\section{Related Work}
\paragraph{Evolving time-series agents.} TimeSage-EV evaluates repeated analysis under changing release periods and strict evidence cutoffs \citep{yao2026timesage}. Its TimeSage-1.0 harness adds reusable skills that become available to later periods. TimeSage-MT provides multi-turn tasks with explicit analytical skill requirements and verifiable task records \citep{kong2026timesagemt}. We use TimeSage-EV to define the longitudinal scientific object and TimeSage-MT to isolate intervention surfaces under matched controls.

\paragraph{Experience-derived agent memory.} Agent Workflow Memory induces reusable procedures from prior trajectories \citep{wang2024awm}. ReasoningBank stores reasoning memories derived from experience \citep{ouyang2026reasoningbank}. MUSE-Autoskill studies skill creation together with later management and evaluation \citep{lin2026muse}. These systems establish persistent procedural state as a useful design pattern. Our question concerns causal credit: which part of a stored skill changes later behavior after generic assistance is matched?

\paragraph{Validated skill evolution.} SkillCommit emphasizes behaviorally validated scope expansion before a skill is committed to broader reuse \citep{he2026skillcommit}. Work on dynamic skill libraries similarly highlights the need to evaluate skills after creation \citep{li2026dynamicskills}. Our protocol adds a matched implementation swap. The visible tool interface stays fixed while the hidden executable procedure changes. This design separates interface prompting from implementation value.

\paragraph{Skill libraries as intervention surfaces.} A modular skill library exposes a causal boundary that ordinary conversation memory does not. The same model can receive the same dialogue while one registered procedure changes. F3 operationalizes this boundary with a native tool loop. The TimeSage chronological contract then extends the same idea across release periods, where a skill written earlier becomes state for a later endpoint.
